\section{Décomposition cyclique et quotients de déterminants}
Les projecteurs cycliques
\[
P_3=\tfrac14(I+S^3+S^6+S^9),\qquad
P_4=\tfrac13(I+S^4+S^8)
\]
ont des images de dimensions trois et quatre. Si $E=P_3-P_4$, alors
\begin{equation}
\frac{\Tr E}{12}=-\frac1{12}.
\label{eq:minus12}
\end{equation}
Le scalaire provient d'une différence de rangs normalisée, non d'une
somme ordinaire d'entiers positifs.

\begin{theorem}[Représentation du lecteur de canal par un quotient de déterminants]
Soient $\mathcal H_3=\operatorname{Fix}(S^3)$ et
$\mathcal H_4=\operatorname{Fix}(S^4)$. Pour $0<s<1$,
\begin{equation}
F(s)=\frac{\det_{\mathcal H_3}(I-sS)}{\det_{\mathcal H_4}(I-sS)}
=\frac{1-s^3}{1-s^4}
=\frac{1+s+s^2}{1+s+s^2+s^3}.
\label{eq:det}
\end{equation}
Son prolongement continu en $s=1$ vaut $3/4$.
\end{theorem}
\begin{proof}
La restriction de $S$ à chaque image est un cycle à trois ou quatre
positions, respectivement. Leurs polynômes caractéristiques donnent
$\det(I-sS)=1-s^r$ pour $r=3,4$. Après simplification du facteur commun $1-s$,
on obtient la dernière expression, régulière en $s=1$.
\end{proof}

La substitution $s=q^{30}$ donne
$F(q^{30})=(1-q^{90})/(1-q^{120})$.
La même paire d’espaces qui produit \eqref{eq:minus12}
produit le quotient de canal. Cette égalité n’identifie pas
une matrice à douze coordonnées à tout l’état TPK.

\subsection{Information orientée et secteur complexe}
Dans le secteur antipodal $V_-=\operatorname{Ran}((I-S^6)/2)$,
l'opérateur $J=S^3$ satisfait $J^2=-I$ et $J^{\mathsf T}=-J$.
Avec $Q=P_4(I-S^6)/2$, on a
\[
\operatorname{rank}Q=2,\qquad E|_{V_-}=-Q,\qquad B|_{V_-}=5I-3Q.
\]
Sur $\operatorname{Ran}Q$, $S^2=-I$ et le déterminant est $1+s^2$.
Ainsi,
\[
F(s)=\frac{1+s+s^2}{1+s}\,\frac1{1+s^2}.
\]
L'inversion d'orientation change le signe de la forme
$\omega(u,v)=\langle Ju,v\rangle$, mais conserve $B$ et $F$.
L'évaluation scalaire ne distingue pas tous les transports orientés.

La trace transportée produit en outre
\[
a_n^{\rm tr}=\Tr(S^nE)=3\mathbf1_{3\mid n}-4\mathbf1_{4\mid n},
\]
et, pour $|s|<1$,
\[
-\log F(s)=\sum_{n\ge1}\frac{a_n^{\rm tr}}n\,s^n.
\]
Cette suite de moments relie la décomposition cyclique à la
représentation spectrale du lecteur. La nature arithmétique de valeurs
particulières, comme la constante de Catalan, nécessite en outre l’analyse
arithmétique correspondante.
